\scanpage{18}
\mspage{16}
ant \deleted{en particulier} à tout $(U\to T)$ le faisceau $\underline{O}_T$ sur $T$, on trouve un faisceau d'anneaux $\underline{O}_C$ sur $C$. Ceci posé, les cristaux de modules (quasi-cohérents) sur $S$ s'identifient aux Modules quasi-cohérents (i.e. localement conoyau d'un homomorphisme de Modules libres) sur $\underline{O}_C$; les cristaux ``de \deleted{type} \inserted{prés.} finie'' i.e. les cristaux de Modules de présentation finie, correspondent exactement aux Modules de présentation finie sur $\underline{O}_C$.

Quand on a un $R$-morphisme $S'\longrightarrow S$, je ne vois pas de morphisme naturel correspondant entre les \uline{sites} cristallogènes correspondants à $S,S'$. Ce n'est pas bien gênant, car introduisant les \uline{topos cristallogènes} \inserted{$\operatorname{Topcr}_{S/R}$ et $\operatorname{Topcr}_{S'/R}$} définis par les sites en question, on définit par la méthode esquissée dans 1.3. un \deleted{homo}morphisme
\[
\operatorname{Topcr}_{S'/R}\longrightarrow\operatorname{Topcr}_{S/R},
\]
correspondant à la notion naturelle de ``image inverse de faisceaux'' \inserted{en sens inverse} (qui, j'avoue, devrait être définie avec le plus grand soin).

1.17. On est donc dans une situation où on peut faire marcher la machine cohomologique. Etant loin de mes bases, je n'ai pas tenté sérieusement de tirer au clair si les images directes supérieures qu'on définit ainsi, et en particulier les groupes de cohomologie du site cristallog\deleted{\unsure{raphique}}\inserted{ène}, donnent ce qu'on voudrait. Le test-clef est le suivant: \uline{si} $R$ \uline{est le spectre d'un corps}, \uline{et si} $S$ \uline{est lisse sur} $R$, \uline{et propre sur} $R$ \uline{dans le cas de la caractéristique} $p>0$, \uline{on voudrait trouver}, \uline{comme cohomologie à coefficients dans} $\underline{O}_C$ \uline{lui-même}, \uline{la cohomologie de De Rham de} $S$ \uline{relativement à} $R$.
\margin{\uline{marche en car. $0$, et \uuline{pas} en car $p>0$} (\unsure{faut|faute} d'introduire \unsure{les} puiss. div.)}
Plus généralement, sans condition sur $R$, si $f:S'\to S$ est propre et lisse, on voudrait trouver comme \inserted{``valeur en $S'$'' de} $R^if_{\mathrm{cris}\,*}(\underline{O}_{C'})$ \inserted{(cf 1.4)} le faisceau de cohomologie de De Rham $\underline{R}^if_*(\underline{\Omega}^{*}_{S'/S})$, \deleted{\unsure{et et comme connexion}} et \inserted{on voudrait} que la connexion canonique à la Gauss-Manin sur ce dernier soit celle associée à la stratification canonique provenant de la structure cristalline (cf 1.8.). L'existence d'une connexion de Gauss-Manin en cohomologie de De-Rham (que j'ai vérifié pour n'importe quel morphisme lisse), et le tapis de Washnitzer-Monsky, sont en tous cas des indications assez fortes pour conduire à penser qu'il doive exister une théorie cohomologique, en termes de cristaux, donnant de telles valeurs pour l'argument $\underline{O}_C$. Au cas où la cohomologie du site cris-
\margin{17}
